\scanpage{19}
\mspage{18}

la chose essentielle est la connaissance de
\[
\underset{\unsure{M|\mathcal{M}}}{\between}\quad\text{\unsure{relt\textsuperscript{n}|relat\textsuperscript{n}|rel\textsuperscript{n}} binaire sym. réfl. (sur }M\text{)}
\tag*{(1.35)}
\]
\deleted{\unsure{qui indique|qui implique}} satisfaisant
\[
\unsure{X|x}\between\unsure{Y|y},\quad X'\leq X,\,Y'\leq Y\quad\Longrightarrow\quad X'\between Y'
\tag*{(1.36)}
\]

\deleted{\unsure{elle|alors}}

il y a corr. biuniv. \unsure{entre|en} l'ens \unsure{des|de} telles relations sur $M$, et l'ens des $\unsure{\mathcal{F}|\mathfrak{F}|\mathscr{F}}\inserted{\subset\unsure{\mathcal{P}|\mathrm{P}|\mathfrak{P}|\mathscr{P}}(M)}$ satisfaisant les conditions énumérées précédemment, plus celle\deleted{s d'être} que les $F\in\mathcal{F}$ soient \deleted{\unsure{toutes|toute|d'être}} de type fini\unsure{.|e|s} C'est en ce sens que j'ai pu dire que ``moralement'', le couple $(M,\mathcal{F}\subset\mathcal{P}(M))$ équivalait à \deleted{\illegible{1}} la donnée de $\underset{M}{\leqslant}$ seul \unsure{de|des} $\between$ satisfaisant (1.36).

\deleted{Mais il s'agit maintenant}

Pratiquement, dans presque tous les exemples, c'est via $M$\inserted{(\unsure{multistrates|multistrate})}, muni de $\leq$ et d'une relation de disjonction $\between$, qu'on d\unsure{é|e}finir\unsure{a|.} les ateliers\ldots

Il faut à présent en venir à la relation $\ll$. \deleted{\unsure{Comme}} Elle va pour l'instant servir surtout à former des ombres, Je commence par les grandes ombres
